% chapters/ch07-validation.tex
\chapter{Design Validation Without SQL}

\begin{learningobjectives}
In this chapter you will learn how to validate a relational design before implementing it.
You will be able to validate a schema by questions, by rule-to-schema traceability, and by controlled test data.
You will also learn a practical review checklist that catches most modelling errors early.
\end{learningobjectives}

\section{Validation is part of design}
A design is not correct because it looks reasonable.
It is correct because it preserves meaning under use.
Validation is the stage where you try to break the model on purpose, using the rules and the questions the system must support.

This book delays SQL, but validation does not require SQL.
Validation is about logic, not syntax.
A schema can be validated with careful reading, traceability, and small thought experiments.

\begin{definitionbox}
\term{Design validation} is the process of checking that a schema:
\begin{itemize}
  \item represents the intended domain meaning,
  \item supports the required questions,
  \item and prevents (or makes detectable) invalid states.
\end{itemize}
\end{definitionbox}

\section{Validation artifacts as deliverables}
Validation should leave artifacts you can review and reuse.
These artifacts make the design auditable and easier to improve.

\begin{figure}[htbp]
\centering
\includegraphics[width=0.85\textwidth]{validation-artifacts.pdf}
\caption{Validation artifacts connect rules to evidence and review notes.}
\end{figure}

\subsection{Questions list}
Keep a short list of key questions that the system must answer.
These questions drive the structure and keep validation grounded in use.

\subsection{Traceability map}
Write a rule-to-schema table that links each rule to constraints or structures.
This makes omissions visible.

\subsection{Paper datasets}
Create small, explicit datasets that include valid and invalid scenarios.
They expose hidden ambiguity before implementation.

\subsection{Review notes}
Capture the decisions and open issues from validation.
These notes become the starting point for the next review.

\section{Method 1: Validate with key questions}
Every domain has questions it must answer.
If your schema cannot express the answers cleanly, the design is not finished.

The point is not to list every possible query.
The point is to identify a small set of \emph{key questions} that reflect the purpose of the system.
These questions should include:
questions about current state, historical state, and exceptional states.

A useful discipline is to write at least:
\begin{itemize}[leftmargin=*,nosep]
  \item 3 questions about lists and lookups,
  \item 3 questions about summaries and counts,
  \item 2 questions about detecting violations or anomalies.
\end{itemize}

\begin{principlebox}
\textbf{Questions reveal structure}\\
If a key question requires awkward joins or duplicated facts, the structure is wrong.
Validate the schema by the questions it must answer.
\end{principlebox}

\subsection{Clinic: key questions}
For the clinic domain, a minimal set could be:
\begin{itemize}
  \item Q1: What appointments does a patient have this week?
  \item Q2: What is a doctor’s schedule for a given day?
  \item Q3: Which patients have no future appointments booked?
  \item Q4: How many appointments did each doctor have last month?
  \item Q5: Are there any double bookings for a doctor?
\end{itemize}

Even without SQL, you can check if each question maps naturally to relations and relationships.
If a question requires duplicated data or unnatural structures, the schema is signalling a problem.

\subsection{Library: key questions}
For the library domain:
\begin{itemize}
  \item Q1: Which copies are available for a given title?
  \item Q2: What active loans does a member have?
  \item Q3: What is the loan history of a copy?
  \item Q4: Which loans are overdue today?
  \item Q5: Are there copies with more than one active loan?
  \item Q6: What is the reservation queue for a title?
\end{itemize}

Notice how these questions force the “title vs copy” separation.
If you cannot answer Q1 and Q3 cleanly, your model is usually collapsing those concepts.

\begin{examplebox}
\textbf{University: key questions}\par\smallskip
\begin{itemize}
  \item Q1: Which courses is a student enrolled in this term?
  \item Q2: Which students are on the waitlist for a course?
  \item Q3: Which courses have no enrolled students yet?
  \item Q4: Are there students enrolled without meeting prerequisites?
\end{itemize}
\end{examplebox}

\section{Method 2: Rule-to-schema traceability}
Rules are the strongest validation tool because they are explicit.
Traceability means that every important rule is linked to the design element that enforces it or supports it.

A traceability map is a short table.
It turns vague assurance (“yes, we modelled that”) into concrete evidence.

\begin{principlebox}
\textbf{Traceability discipline}\\[2pt]
For each important rule:
\begin{itemize}[leftmargin=*,nosep]
  \item identify where it is enforced (constraint), or
  \item identify how it is supported (structure), or
  \item explicitly state that enforcement is external and how violations are detected.
\end{itemize}
\end{principlebox}

\begin{examplebox}
\textbf{Clinic traceability example}\\[4pt]
\begin{tabularx}{\textwidth}{@{}lX@{}}
\toprule
Rule & Design element(s) \\
\midrule
Appointment has exactly one patient & Appointment.patient\_id is mandatory FK to Patient \\
Appointment has exactly one doctor & Appointment.doctor\_id is mandatory FK to Doctor \\
No double booking (slot-based) & Unique(doctor\_id, start\_time) on Appointment \\
\bottomrule
\end{tabularx}
\end{examplebox}

\begin{examplebox}
\textbf{Library traceability example}\\[4pt]
\begin{tabularx}{\textwidth}{@{}lX@{}}
\toprule
Rule & Design element(s) \\
\midrule
Each Copy belongs to one Title & Copy.title\_id is mandatory FK to Title \\
Each Loan references one Member & Loan.member\_id is mandatory FK to Member \\
Each Loan references one Copy & Loan.copy\_id is mandatory FK to Copy \\
Due date not before checkout date & Loan.due\_date \(\ge\) Loan.checkout\_date \\
One active loan per copy & No two Loans with same copy\_id and missing return\_date \\
Reservation references a Title & Reservation.title\_id is mandatory FK to Title \\
Reservation references a Member & Reservation.member\_id is mandatory FK to Member \\
Copy barcode unique when present & Unique(Copy.barcode) when not null \\
Loan history preserved & Loan is an event relation, not a status field on Copy \\
\bottomrule
\end{tabularx}
\end{examplebox}

Traceability also reveals missing design elements.
If you have a rule and cannot point to any place in the design that supports it, the model is incomplete.

\section{Method 3: Validate with controlled test data (on paper)}
Even without SQL, you can validate by constructing small example datasets.
This is an underrated technique.
You propose a handful of rows for each relation and see whether:
\begin{itemize}[leftmargin=*,nosep]
  \item valid scenarios are representable without awkward duplication,
  \item invalid scenarios are prevented or at least visible as violations,
  \item history can be represented when time is involved.
\end{itemize}

This method is especially useful for business constraints.

\subsection{Clinic thought experiment}
Try to represent:
\begin{itemize}
  \item a doctor with no appointments yet,
  \item a patient with two future appointments,
  \item an invalid double booking attempt.
\end{itemize}

If your design cannot represent the first two scenarios, it is too restrictive.
If it cannot represent the invalid scenario as an explicit violation, it is too permissive.

\begin{examplebox}
\textbf{Clinic paper dataset (sketch)}\par\smallskip
\begin{tabular}{@{}lll@{}}
\toprule
\multicolumn{3}{@{}l@{}}{\textbf{Patient}} \\
\midrule
patient\_id & name & birth\_date \\
P1 & Ana Costa & 1990-04-11 \\
P2 & Rui Silva & 1985-09-03 \\
P3 & Marta Lima & 2001-01-22 \\
\bottomrule
\end{tabular}
\par\smallskip
\begin{tabular}{@{}lll@{}}
\toprule
\multicolumn{3}{@{}l@{}}{\textbf{Doctor}} \\
\midrule
doctor\_id & name & specialty \\
D1 & Joana Reis & cardiology \\
D2 & Hugo Pereira & pediatrics \\
D3 & Sara Pinto & dermatology \\
\bottomrule
\end{tabular}
\par\smallskip
\begin{tabular}{@{}llll@{}}
\toprule
\multicolumn{4}{@{}l@{}}{\textbf{Appointment}} \\
\midrule
appointment\_id & start\_time & patient\_id & doctor\_id \\
A1 & 2026-02-03 09:00 & P1 & D1 \\
A2 & 2026-02-03 09:30 & P2 & D1 \\
A3 & 2026-02-03 10:00 & P3 & D2 \\
\bottomrule
\end{tabular}
\end{examplebox}

\subsection{Library thought experiment}
Try to represent:
\begin{itemize}
  \item one title with two copies,
  \item one copy loaned and returned,
  \item one copy loaned to two members at the same time (invalid),
  \item two reservations for the same title.
\end{itemize}

A good design makes the invalid state either impossible (ideal) or immediately detectable.
If the invalid state blends in without obvious contradiction, you need stronger constraints or clearer state modelling.

\begin{examplebox}
\textbf{Library test rows (sketch)}\par\smallskip
\[
\rel{Copy}(\underline{copy\_id}, status, title\_id \rightarrow \rel{Title})
\]
Rows: (C1, on\_loan, T1), (C2, available, T1)\par
\[
\rel{Loan}(\underline{loan\_id}, copy\_id \rightarrow \rel{Copy}, member\_id \rightarrow \rel{Member}, return\_date?)
\]
Rows: (L1, C1, M1, \text{missing}), (L2, C1, M2, \text{missing})\par
If both L1 and L2 are allowed, the “one active loan per copy” rule is violated.
\end{examplebox}

\begin{notebox}
Keep test datasets small and explicit.
Two or three rows per relation are usually enough to expose hidden ambiguity.
\end{notebox}

\section{Common mistakes}
Most beginner designs fail in repeatable ways.
Validation methods catch them early.

\subsection{Collapsing distinct concepts}
Title vs copy, customer vs address, order vs payment, doctor vs specialty.
If you collapse them, you will struggle to represent history, multiplicity, or policy.
Key questions usually reveal this immediately.

\subsection{Treating events as static relationships}
Loans, appointments, payments, enrollments, subscriptions.
These are events with time and state.
If you model them as a single foreign key or a simple relationship, you lose history and constraints.
Controlled test data exposes the problem fast.

\subsection{Missing optionality decisions}
If you do not decide what is mandatory, your schema quietly allows invalid states.
Traceability catches this because rules like “exactly one” demand a mandatory reference.

\subsection{Encoding collections inside attributes}
Comma-separated lists, repeated columns, and JSON blobs used as a shortcut.
These break 1NF and make constraints hard to express.
Controlled test data exposes the awkwardness, and normalization principles give the fix.

\section{Limits of validation}
Validation improves correctness, but it cannot guarantee it.
Some errors only appear under real data volume or evolving business rules.

\begin{notebox}
\textbf{What validation cannot guarantee}\\
Validation cannot prove a design is perfect.
It can only reduce risk by exposing contradictions, missing rules, and weak assumptions.
\end{notebox}

\section{A practical review checklist}
Use this checklist to review any schema before implementation.

\begin{examplebox}
\textbf{Schema review checklist}\par\smallskip
\begin{itemize}
  \item Identity: Does every entity have a clear primary key? Is it stable?
  \item Relationships: Are 1:N and N:M represented correctly (FK vs association relation)?
  \item Optionality: For each FK, is it mandatory or optional, and does that match the rules?
  \item Redundancy: Are you repeating the same facts across multiple relations?
  \item Time: Are time-based events represented as entities (Loan, Appointment) with their own attributes?
  \item Constraints: Are key business constraints stated and mapped to enforcement or detection?
  \item Questions: Can you answer the key questions naturally from the structure?
\end{itemize}
\end{examplebox}

This checklist is short on purpose.
If you apply it consistently, it catches most structural issues before a single line of SQL exists.

\section{Mini checklist}
Before moving to implementation, confirm:
\begin{itemize}[leftmargin=*,nosep]
  \item Each key question has a clear path through the schema.
  \item Every business rule is mapped to a constraint or detection plan.
  \item Test rows can represent both valid and invalid scenarios.
\end{itemize}

\begin{chaptersummary}
Validation is design work.
You can validate a schema without SQL by checking whether it supports key questions, by building a rule-to-schema traceability map, and by constructing small controlled datasets to test representability and constraint behaviour.
The goal is to ensure your model preserves meaning under real use, not only on paper.
\end{chaptersummary}

\begin{exercisebox}
Pick one domain (clinic, library, university, online shop) and answer in full sentences.
\begin{enumerate}[label=\arabic*.]
  \item Write 8 key questions (at least 2 about violations or anomalies).
  \item Write 10 business rules and create a traceability table for at least 6 of them.
  \item Propose a small dataset (3--5 rows per relation) that represents two valid scenarios.
  \item Propose one invalid scenario and explain how your schema would prevent or detect it.
  \item Use the checklist to find one improvement to your current model.
  \item Identify one common mistake your schema might still contain and explain how you would test for it.
  \item Add one question that forces a decision about optionality and explain how the schema answers it.
  \item Write a short validation note that documents one assumption and one open issue.
  \item Create a traceability row for a constraint that cannot be enforced directly and state how to detect violations.
\end{enumerate}
\end{exercisebox}

In the next chapter we focus on naming and documentation.
The artifacts you created here become inputs to those decisions.
